\chapter*{Conclusion}
\label{chap:conclusion}
\addstarredchapter{Conclusion}

\section*{Rappel du problème}
Parvenu au terme de notre travail, il était question pour nous d'effectuer le déploiement et le passage à l'échelle d'une application
multi-tiers. L'application étant web, disposant d'une API faite en JavaScript et d'une base de données géospatiale, ceci indiquant qu'il 
s'agisse d'un SIG où le temps de calcul est important.

\section*{Démarche et résultats}
Nous avons choisi pour résoudre notre problème, une approche itérative incrémentale où on avançait étape par étape et où on améliorait ce 
qui avait déjà été fait. Pour y arriver, on a :
\begin{itemize}
    \item Appris à nous connecter à l'ordinateur distant,
    \item Appris à utiliser les outils nécessaires,
    \item Effectué une conteneurisation progressive,
    \item \'Ecrit un générateur de charge et effectué des tests de charge,
    \item Procédé à un déploiement sur Kubernetes.
\end{itemize}

\section*{Limites}
L'on déplore :
\begin{itemize}
    \item Le temps de chargement de la base de données lors de la première initialisation,
    \item Le temps d'initialisation du cluster qui est plus grand en fonction du nombre de n\oe ud à cause de la base de données.
\end{itemize}

\section*{Perpective}
On peut :
\begin{itemize}
    \item Mettre chacune des tables la base de données dans un conteneur Docker, le tout répartis dans un pod Kubernetes pour diviser le temps d'initialisation par 2 pour un n\oe ud.
\end{itemize}